% --- Archon physics formalization source begin ---
% archon:physics
% archon:covers PhyXMiniProblems/problem_phyx_mini_0400.lean
% archon:source-report reports/phyx_mini/problem_phyx_mini_0400.source.json

\chapter{Physics problem phyx\_mini\_0400}
\label{ch:PhyXMiniProblems_problem_phyx_mini_0400}

\paragraph{Problem source.}
\#\# Physical scenario

A 1-m\$\textasciicircum{}3\$ rigid tank has propane at 100 kPa, 300 K and connected by a valve to another tank of 0.5 m\$\textasciicircum{}3\$ with propane at 250 kPa, 400 K. The valve is opened, and the two tanks come to a uniform state at 325 K.

\#\# Question

What is the final pressure?

\#\# Answer choices

- A: A: 4MPa

- B: B: 139.9kPa

- C: C: 154kPa

- D: D: 10.2kPa

\#\# Figure caption (auxiliary; use the image as primary evidence)

The image depicts a schematic diagram of two tanks connected by pipes. 

- **Tanks**: There are two cylindrical tanks labeled "A" and "B." Tank A is positioned above tank B.
- **Pipes**: Pipes connect the two tanks with a flow path that goes from the bottom of tank A to the top of tank B.
- **Valve**: There is a valve depicted on the pipe between the two tanks, indicating a control point in the flow path.
- **Flow Direction**: The design suggests a fluid can move from tank A to tank B, possibly controlled by the valve. 
- **Color**: The tanks are shaded in blue, with tank A being lighter and tank B being darker.

The overall setup implies a fluid transfer or mixing system where the flow can be regulated.

\#\# Dataset metadata

Ideal Gases and Kinetic Theory; Multi-Formula Reasoning, Physical Model Grounding Reasoning

\paragraph{Recorded answer/context.}
B: B: 139.9kPa

\paragraph{Figure/image path.}
/root/proposal\_for\_physic/science-mango/phyx\_mini\_run/phyx\_data/test\_image/400.png

\paragraph{Formalization target.}
create a compiling Lean file with sorry bodies at `PhyXMiniProblems/problem\_phyx\_mini\_0400.lean`. The Lean declarations must preserve the physical quantities, dimensions or dimensional roles, figure labels, governing-law hypotheses, and final relation expressed by this problem.
Use Mathlib/Physlib names found through LeanExplore where available. If a physics API is missing, introduce faithful local abstractions rather than scalar placeholder aliases.

\begin{theorem}[Physics formalization target]
\label{thm:physics:phyx_mini_0400:target}
\lean{PhyXMiniProblems.ProblemPhyXMini0400.finalPressureInKilopascals_exact_and_matches_recordedAnswerB}
\uses{def:physics:phyx-mini-0400:phyxminiproblems-problemphyxmini0400-pressureinkilopascals, def:physics:phyx-mini-0400:phyxminiproblems-problemphyxmini0400-connectedpropanetanks, def:physics:phyx-mini-0400:phyxminiproblems-problemphyxmini0400-matchesprimaryfigure, def:physics:phyx-mini-0400:phyxminiproblems-problemphyxmini0400-matchesproblemstatement, def:physics:phyx-mini-0400:phyxminiproblems-problemphyxmini0400-hasphysicalparameters, def:physics:phyx-mini-0400:phyxminiproblems-problemphyxmini0400-obeysidealgasandamountconservation, def:physics:phyx-mini-0400:phyxminiproblems-problemphyxmini0400-reachesuniformpressure, def:physics:phyx-mini-0400:phyxminiproblems-problemphyxmini0400-finalpressureinkilopascals, def:physics:phyx-mini-0400:phyxminiproblems-problemphyxmini0400-recordedanswerchoice, def:physics:phyx-mini-0400:phyxminiproblems-problemphyxmini0400-agreeswithonedecimalkilopascalrounding}
The assigned autoformalize agent should translate this physics problem into Lean declarations in the covered file, with theorem and lemma proof bodies written as `by sorry`.
\end{theorem}
\begin{proof}
This is an autoformalization task, not a proof task. Produce statements that can later be proved by the physics prover without weakening the physical model.
\end{proof}

% --- Archon declaration topology begin ---
\section{Declaration topology}
\begin{definition}[Gas Volume]
  \label{def:physics:phyx-mini-0400:phyxminiproblems-problemphyxmini0400-gasvolume}
  \lean{PhyXMiniProblems.ProblemPhyXMini0400.GasVolume}
  A physical gas volume, carrying the length-cubed dimension
\end{definition}
\begin{proof}
  This is the declared model definition of Gas Volume.
\end{proof}
\begin{definition}[Molar Gas Constant Quantity]
  \label{def:physics:phyx-mini-0400:phyxminiproblems-problemphyxmini0400-molargasconstantquantity}
  \lean{PhyXMiniProblems.ProblemPhyXMini0400.MolarGasConstantQuantity}
  The molar gas constant with its energy-per-temperature dimension. Physlib's Dimension has no amount-of-substance coordinate. Its inverse-mole role is therefore recorded by the type name and by the explicit mole readouts paired with this quantity in the ideal-gas law below
\end{definition}
\begin{proof}
  This is the declared model definition of Molar Gas Constant Quantity.
\end{proof}
\begin{definition}[volume In Cubic Meters]
  \label{def:physics:phyx-mini-0400:phyxminiproblems-problemphyxmini0400-volumeincubicmeters}
  \lean{PhyXMiniProblems.ProblemPhyXMini0400.volumeInCubicMeters}
  \uses{def:physics:phyx-mini-0400:phyxminiproblems-problemphyxmini0400-gasvolume}
  Read a physical gas volume in coherent SI cubic metres
\end{definition}
\begin{proof}
  This is the declared model definition of volume In Cubic Meters.
\end{proof}
\begin{definition}[pressure In Pascals]
  \label{def:physics:phyx-mini-0400:phyxminiproblems-problemphyxmini0400-pressureinpascals}
  \lean{PhyXMiniProblems.ProblemPhyXMini0400.pressureInPascals}
  Read a physical pressure in coherent SI pascals
\end{definition}
\begin{proof}
  This is the declared model definition of pressure In Pascals.
\end{proof}
\begin{definition}[pressure In Kilopascals]
  \label{def:physics:phyx-mini-0400:phyxminiproblems-problemphyxmini0400-pressureinkilopascals}
  \lean{PhyXMiniProblems.ProblemPhyXMini0400.pressureInKilopascals}
  \uses{def:physics:phyx-mini-0400:phyxminiproblems-problemphyxmini0400-pressureinpascals}
  Read a physical pressure in kilopascals
\end{definition}
\begin{proof}
  This is the declared model definition of pressure In Kilopascals.
\end{proof}
\begin{definition}[temperature In Kelvins]
  \label{def:physics:phyx-mini-0400:phyxminiproblems-problemphyxmini0400-temperatureinkelvins}
  \lean{PhyXMiniProblems.ProblemPhyXMini0400.temperatureInKelvins}
  Read an absolute Temperature in kelvins. The storage scale is explicit because Physlib's Temperature permits an arbitrary zero-preserving unit
\end{definition}
\begin{proof}
  This is the declared model definition of temperature In Kelvins.
\end{proof}
\begin{definition}[molar Gas Constant In Joules Per Mole Kelvin]
  \label{def:physics:phyx-mini-0400:phyxminiproblems-problemphyxmini0400-molargasconstantinjoulespermolekelvin}
  \lean{PhyXMiniProblems.ProblemPhyXMini0400.molarGasConstantInJoulesPerMoleKelvin}
  \uses{def:physics:phyx-mini-0400:phyxminiproblems-problemphyxmini0400-molargasconstantquantity}
  Read the molar gas constant in joules per mole-kelvin
\end{definition}
\begin{proof}
  This is the declared model definition of molar Gas Constant In Joules Per Mole Kelvin.
\end{proof}
\begin{definition}[Tank Label]
  \label{def:physics:phyx-mini-0400:phyxminiproblems-problemphyxmini0400-tanklabel}
  \lean{PhyXMiniProblems.ProblemPhyXMini0400.TankLabel}
  The tank labels printed in the primary figure
\end{definition}
\begin{proof}
  This is the declared model definition of Tank Label.
\end{proof}
\begin{definition}[Process Stage]
  \label{def:physics:phyx-mini-0400:phyxminiproblems-problemphyxmini0400-processstage}
  \lean{PhyXMiniProblems.ProblemPhyXMini0400.ProcessStage}
  The initially separated state and the final uniform state
\end{definition}
\begin{proof}
  This is the declared model definition of Process Stage.
\end{proof}
\begin{definition}[Vertical Position]
  \label{def:physics:phyx-mini-0400:phyxminiproblems-problemphyxmini0400-verticalposition}
  \lean{PhyXMiniProblems.ProblemPhyXMini0400.VerticalPosition}
  Vertical placement of a tank in the supplied schematic
\end{definition}
\begin{proof}
  This is the declared model definition of Vertical Position.
\end{proof}
\begin{definition}[Tank Shape]
  \label{def:physics:phyx-mini-0400:phyxminiproblems-problemphyxmini0400-tankshape}
  \lean{PhyXMiniProblems.ProblemPhyXMini0400.TankShape}
  Shape in which each rigid tank is drawn
\end{definition}
\begin{proof}
  This is the declared model definition of Tank Shape.
\end{proof}
\begin{definition}[Pipe Attachment]
  \label{def:physics:phyx-mini-0400:phyxminiproblems-problemphyxmini0400-pipeattachment}
  \lean{PhyXMiniProblems.ProblemPhyXMini0400.PipeAttachment}
  The attachment location of the connecting pipe at a tank
\end{definition}
\begin{proof}
  This is the declared model definition of Pipe Attachment.
\end{proof}
\begin{definition}[Tank Shade]
  \label{def:physics:phyx-mini-0400:phyxminiproblems-problemphyxmini0400-tankshade}
  \lean{PhyXMiniProblems.ProblemPhyXMini0400.TankShade}
  The blue fill shades distinguishing the two tanks in the image
\end{definition}
\begin{proof}
  This is the declared model definition of Tank Shade.
\end{proof}
\begin{definition}[Valve Position]
  \label{def:physics:phyx-mini-0400:phyxminiproblems-problemphyxmini0400-valveposition}
  \lean{PhyXMiniProblems.ProblemPhyXMini0400.ValvePosition}
  Position of the control valve at a modeled stage
\end{definition}
\begin{proof}
  This is the declared model definition of Valve Position.
\end{proof}
\begin{definition}[Tank Mechanical Model]
  \label{def:physics:phyx-mini-0400:phyxminiproblems-problemphyxmini0400-tankmechanicalmodel}
  \lean{PhyXMiniProblems.ProblemPhyXMini0400.TankMechanicalModel}
  Mechanical model stated for both vessels
\end{definition}
\begin{proof}
  This is the declared model definition of Tank Mechanical Model.
\end{proof}
\begin{definition}[Connection Model]
  \label{def:physics:phyx-mini-0400:phyxminiproblems-problemphyxmini0400-connectionmodel}
  \lean{PhyXMiniProblems.ProblemPhyXMini0400.ConnectionModel}
  Idealization of the pipe shown between the two tanks
\end{definition}
\begin{proof}
  This is the declared model definition of Connection Model.
\end{proof}
\begin{definition}[Gas Model]
  \label{def:physics:phyx-mini-0400:phyxminiproblems-problemphyxmini0400-gasmodel}
  \lean{PhyXMiniProblems.ProblemPhyXMini0400.GasModel}
  Equation-of-state model used for the propane
\end{definition}
\begin{proof}
  This is the declared model definition of Gas Model.
\end{proof}
\begin{definition}[Gas Species]
  \label{def:physics:phyx-mini-0400:phyxminiproblems-problemphyxmini0400-gasspecies}
  \lean{PhyXMiniProblems.ProblemPhyXMini0400.GasSpecies}
  Gas species distinguished by the problem statement
\end{definition}
\begin{proof}
  This is the declared model definition of Gas Species.
\end{proof}
\begin{definition}[Gas Portion]
  \label{def:physics:phyx-mini-0400:phyxminiproblems-problemphyxmini0400-gasportion}
  \lean{PhyXMiniProblems.ProblemPhyXMini0400.GasPortion}
  \uses{def:physics:phyx-mini-0400:phyxminiproblems-problemphyxmini0400-gasspecies}
  A physical gas portion, retaining both its species and its calibrated amount readout. The portion itself is not identified with a bare real number
\end{definition}
\begin{proof}
  This is the declared model definition of Gas Portion.
\end{proof}
\begin{definition}[Two Tank Valve Figure]
  \label{def:physics:phyx-mini-0400:phyxminiproblems-problemphyxmini0400-twotankvalvefigure}
  \lean{PhyXMiniProblems.ProblemPhyXMini0400.TwoTankValveFigure}
  \uses{def:physics:phyx-mini-0400:phyxminiproblems-problemphyxmini0400-tanklabel, def:physics:phyx-mini-0400:phyxminiproblems-problemphyxmini0400-verticalposition, def:physics:phyx-mini-0400:phyxminiproblems-problemphyxmini0400-tankshape, def:physics:phyx-mini-0400:phyxminiproblems-problemphyxmini0400-pipeattachment, def:physics:phyx-mini-0400:phyxminiproblems-problemphyxmini0400-tankshade}
  Labels and geometry read directly from the primary image. In particular, the bitmap shows no flow-direction arrow; flow after opening is part of the process model, not an invented graphical readout
\end{definition}
\begin{proof}
  This is the declared model definition of Two Tank Valve Figure.
\end{proof}
\begin{definition}[Connected Propane Tanks]
  \label{def:physics:phyx-mini-0400:phyxminiproblems-problemphyxmini0400-connectedpropanetanks}
  \lean{PhyXMiniProblems.ProblemPhyXMini0400.ConnectedPropaneTanks}
  \uses{def:physics:phyx-mini-0400:phyxminiproblems-problemphyxmini0400-gasvolume, def:physics:phyx-mini-0400:phyxminiproblems-problemphyxmini0400-molargasconstantquantity, def:physics:phyx-mini-0400:phyxminiproblems-problemphyxmini0400-tanklabel, def:physics:phyx-mini-0400:phyxminiproblems-problemphyxmini0400-processstage, def:physics:phyx-mini-0400:phyxminiproblems-problemphyxmini0400-valveposition, def:physics:phyx-mini-0400:phyxminiproblems-problemphyxmini0400-tankmechanicalmodel, def:physics:phyx-mini-0400:phyxminiproblems-problemphyxmini0400-connectionmodel, def:physics:phyx-mini-0400:phyxminiproblems-problemphyxmini0400-gasmodel, def:physics:phyx-mini-0400:phyxminiproblems-problemphyxmini0400-gasportion, def:physics:phyx-mini-0400:phyxminiproblems-problemphyxmini0400-twotankvalvefigure}
  The two tanks, their thermodynamic observables at both stages, their propane portions, and the connecting equipment. The final pressure is an independent observable; no numerical answer is stored or defined here
\end{definition}
\begin{proof}
  This is the declared model definition of Connected Propane Tanks.
\end{proof}
\begin{definition}[Matches Primary Figure]
  \label{def:physics:phyx-mini-0400:phyxminiproblems-problemphyxmini0400-matchesprimaryfigure}
  \lean{PhyXMiniProblems.ProblemPhyXMini0400.MatchesPrimaryFigure}
  \uses{def:physics:phyx-mini-0400:phyxminiproblems-problemphyxmini0400-connectedpropanetanks}
  Primary-image evidence: A is the upper light-blue horizontal cylinder and B the lower dark-blue one. A visible pipe runs from the bottom of A to the top of B through a visible valve on the left
\end{definition}
\begin{proof}
  This is the declared model definition of Matches Primary Figure.
\end{proof}
\begin{definition}[Matches Problem Statement]
  \label{def:physics:phyx-mini-0400:phyxminiproblems-problemphyxmini0400-matchesproblemstatement}
  \lean{PhyXMiniProblems.ProblemPhyXMini0400.MatchesProblemStatement}
  \uses{def:physics:phyx-mini-0400:phyxminiproblems-problemphyxmini0400-volumeincubicmeters, def:physics:phyx-mini-0400:phyxminiproblems-problemphyxmini0400-pressureinkilopascals, def:physics:phyx-mini-0400:phyxminiproblems-problemphyxmini0400-temperatureinkelvins, def:physics:phyx-mini-0400:phyxminiproblems-problemphyxmini0400-connectedpropanetanks}
  Numerical and qualitative data stated in the problem. These fields specify the two initial states and the final temperature, but no numerical final pressure
\end{definition}
\begin{proof}
  This is the declared model definition of Matches Problem Statement.
\end{proof}
\begin{definition}[Has Physical Parameters]
  \label{def:physics:phyx-mini-0400:phyxminiproblems-problemphyxmini0400-hasphysicalparameters}
  \lean{PhyXMiniProblems.ProblemPhyXMini0400.HasPhysicalParameters}
  \uses{def:physics:phyx-mini-0400:phyxminiproblems-problemphyxmini0400-volumeincubicmeters, def:physics:phyx-mini-0400:phyxminiproblems-problemphyxmini0400-pressureinpascals, def:physics:phyx-mini-0400:phyxminiproblems-problemphyxmini0400-temperatureinkelvins, def:physics:phyx-mini-0400:phyxminiproblems-problemphyxmini0400-molargasconstantinjoulespermolekelvin, def:physics:phyx-mini-0400:phyxminiproblems-problemphyxmini0400-connectedpropanetanks}
  Positivity conditions selecting physically meaningful states and data
\end{definition}
\begin{proof}
  This is the declared model definition of Has Physical Parameters.
\end{proof}
\begin{definition}[Obeys Ideal Gas And Amount Conservation]
  \label{def:physics:phyx-mini-0400:phyxminiproblems-problemphyxmini0400-obeysidealgasandamountconservation}
  \lean{PhyXMiniProblems.ProblemPhyXMini0400.ObeysIdealGasAndAmountConservation}
  \uses{def:physics:phyx-mini-0400:phyxminiproblems-problemphyxmini0400-volumeincubicmeters, def:physics:phyx-mini-0400:phyxminiproblems-problemphyxmini0400-pressureinpascals, def:physics:phyx-mini-0400:phyxminiproblems-problemphyxmini0400-temperatureinkelvins, def:physics:phyx-mini-0400:phyxminiproblems-problemphyxmini0400-molargasconstantinjoulespermolekelvin, def:physics:phyx-mini-0400:phyxminiproblems-problemphyxmini0400-connectedpropanetanks}
  Macroscopic ideal-gas and closed-system amount-conservation laws. The SI readout equation is dimensionally coherent because Pa · m³ = J. Neither law contains a numerical final pressure
\end{definition}
\begin{proof}
  This is the declared model definition of Obeys Ideal Gas And Amount Conservation.
\end{proof}
\begin{definition}[Reaches Uniform Pressure]
  \label{def:physics:phyx-mini-0400:phyxminiproblems-problemphyxmini0400-reachesuniformpressure}
  \lean{PhyXMiniProblems.ProblemPhyXMini0400.ReachesUniformPressure}
  \uses{def:physics:phyx-mini-0400:phyxminiproblems-problemphyxmini0400-pressureinpascals, def:physics:phyx-mini-0400:phyxminiproblems-problemphyxmini0400-connectedpropanetanks}
  Opening the valve allows mechanical equilibration, so the final uniform state has one pressure. This is a qualitative equilibrium law, not the requested numerical pressure
\end{definition}
\begin{proof}
  This is the declared model definition of Reaches Uniform Pressure.
\end{proof}
\begin{definition}[final Pressure In Kilopascals]
  \label{def:physics:phyx-mini-0400:phyxminiproblems-problemphyxmini0400-finalpressureinkilopascals}
  \lean{PhyXMiniProblems.ProblemPhyXMini0400.finalPressureInKilopascals}
  \uses{def:physics:phyx-mini-0400:phyxminiproblems-problemphyxmini0400-pressureinkilopascals, def:physics:phyx-mini-0400:phyxminiproblems-problemphyxmini0400-connectedpropanetanks}
  The common final-pressure readout, taken from tank A, in kilopascals
\end{definition}
\begin{proof}
  This is the declared model definition of final Pressure In Kilopascals.
\end{proof}
\begin{definition}[Answer Choice]
  \label{def:physics:phyx-mini-0400:phyxminiproblems-problemphyxmini0400-answerchoice}
  \lean{PhyXMiniProblems.ProblemPhyXMini0400.AnswerChoice}
  Labels of the four displayed answer choices
\end{definition}
\begin{proof}
  This is the declared model definition of Answer Choice.
\end{proof}
\begin{definition}[pressure In Kilopascals]
  \label{def:physics:phyx-mini-0400:phyxminiproblems-problemphyxmini0400-answerchoice-pressureinkilopascals}
  \lean{PhyXMiniProblems.ProblemPhyXMini0400.AnswerChoice.pressureInKilopascals}
  \uses{def:physics:phyx-mini-0400:phyxminiproblems-problemphyxmini0400-pressureinkilopascals, def:physics:phyx-mini-0400:phyxminiproblems-problemphyxmini0400-answerchoice}
  Kilopascal value printed beside each displayed answer label
\end{definition}
\begin{proof}
  This is the declared model definition of pressure In Kilopascals.
\end{proof}
\begin{definition}[recorded Answer Choice]
  \label{def:physics:phyx-mini-0400:phyxminiproblems-problemphyxmini0400-recordedanswerchoice}
  \lean{PhyXMiniProblems.ProblemPhyXMini0400.recordedAnswerChoice}
  \uses{def:physics:phyx-mini-0400:phyxminiproblems-problemphyxmini0400-answerchoice}
  Answer label recorded by the supplied dataset
\end{definition}
\begin{proof}
  This is the declared model definition of recorded Answer Choice.
\end{proof}
\begin{definition}[Agrees With One Decimal Kilopascal Rounding]
  \label{def:physics:phyx-mini-0400:phyxminiproblems-problemphyxmini0400-agreeswithonedecimalkilopascalrounding}
  \lean{PhyXMiniProblems.ProblemPhyXMini0400.AgreesWithOneDecimalKilopascalRounding}
  Agreement with a displayed pressure rounded to the nearest 0.1 kPa: the exact value is strictly within half of one tenth of the displayed value
\end{definition}
\begin{proof}
  This is the declared model definition of Agrees With One Decimal Kilopascal Rounding.
\end{proof}
% --- Archon declaration topology end ---
% --- Archon physics formalization source end ---
